\chapter{Ficheros de resultados}
\label{anexo:resultados}

Este anexo documenta cada uno de los ficheros de resultados que produce
el pipeline y sobre los que se alimenta tanto esta memoria como la
interfaz web. Todos ellos residen en la raíz del repositorio y son la
única fuente de verdad de las cifras del trabajo: cualquier cifra citada
en la memoria puede trazarse a una fila o campo concreto de estos
ficheros.

\section*{Auditoría de cohortes}
\phantomsection
\addcontentsline{toc}{section}{Auditoría de cohortes}

\begin{description}
  \item[\texttt{AUDITORIA\_COHORTES.csv}] Una fila por cohorte
  descargada. Columnas relevantes: \texttt{Cohorte}, \texttt{Plataforma},
  \texttt{N\_Total}, \texttt{N\_Sano}, \texttt{N\_Enfermo},
  \texttt{N\_Sin\_Clasificar}, \texttt{Tasa\_Exito\_Curacion},
  \texttt{Evaluable\_Como\_Test}. Es la base de la tabla de datos del
  capítulo~\ref{cap:metodos} (tabla~\ref{tab:cohortes}) y de todas las
  cifras de cohortes y muestras del hero de la interfaz.
\end{description}

\section*{Rendimiento del modelo (LODO)}
\phantomsection
\addcontentsline{toc}{section}{Rendimiento del modelo (LODO)}

\begin{description}
  \item[\texttt{LODO\_HONESTO\_RESULTADOS.csv}] Una fila por cohorte
  usada como test en LODO para la tarea tumor \emph{vs.} sano. Columnas:
  \texttt{Cohorte\_Test}, \texttt{n\_test}, \texttt{n\_Sano},
  \texttt{n\_Enfermo}, \texttt{n\_train}, \texttt{Evaluable},
  \texttt{Baseline\_Mayoritaria}, \texttt{Accuracy},
  \texttt{Balanced\_Accuracy}, \texttt{AUC}, \texttt{Sensibilidad},
  \texttt{Especificidad}, \texttt{Ganancia\_vs\_Baseline},
  \texttt{Supera\_Baseline}, \texttt{Genes\_Seleccionados}. Alimenta la
  tabla~\ref{tab:lodo_tvs} y las cifras del capítulo~\ref{cap:resultados}
  sección~\ref{sec:res_tvs}.

  \item[\texttt{SUBTIPO\_LODO\_RESULTADOS.csv}] Análogo al anterior pero
  para la tarea ADC \emph{vs.} SQC. Columnas: \texttt{Cohorte\_Test},
  \texttt{n\_test}, \texttt{n\_ADC}, \texttt{n\_SQC}, \texttt{n\_train},
  \texttt{Baseline\_Mayoritaria}, \texttt{Accuracy},
  \texttt{Balanced\_Accuracy}, \texttt{AUC}, \texttt{Sensibilidad\_SQC},
  \texttt{Especificidad\_ADC}, \texttt{Ganancia\_vs\_Baseline},
  \texttt{Genes\_Seleccionados}. Alimenta la tabla~\ref{tab:lodo_sub}.
\end{description}

\section*{Firma génica validada}
\phantomsection
\addcontentsline{toc}{section}{Firma génica validada}

\begin{description}
  \item[\texttt{FIRMA\_VALIDADA\_COMPLETA.csv}] Un gen por fila, con las
  columnas: \texttt{ID\_REF} (símbolo HUGO), \texttt{GSE30219},
  \texttt{GSE50081}, \texttt{GSE19188} (d de Cohen en cada cohorte),
  \texttt{d\_Media}, \texttt{d\_Minima\_Abs}, \texttt{Direccion}. Contiene
  los 1174 genes de la firma consenso multi-cohorte. Es el fichero que
  consulta el buscador de gen del dashboard.

  \item[\texttt{FIRMA\_VALIDADA\_TOP60.csv}] Los 60 primeros genes de la
  firma, con columnas adicionales: \texttt{Rango} (posición en el
  ranking), \texttt{Marcador\_IHC\_Clinica} (booleano),
  \texttt{En\_Panel\_Minimo} (booleano).

  \item[\texttt{FIRMA\_VALIDADA\_RESUMEN.json}] Documento JSON con las
  cifras resumen del análisis:
  \begin{itemize}
    \item \texttt{tarea\_validada}: cadena que identifica la tarea (aquí
    ``Adenocarcinoma frente a carcinoma escamoso'').
    \item \texttt{n\_cohortes\_independientes}: 3.
    \item \texttt{n\_genes\_evaluados}: 22\,880.
    \item \texttt{n\_genes\_validados}: 1174.
    \item \texttt{pct\_genes\_validados}: 5,131.
    \item \texttt{n\_genes\_validados\_tumor\_vs\_sano}: número de genes
    de la firma que también aparecen en la firma tumor-vs-sano.
    \item \texttt{panel\_minimo}: 20.
    \item \texttt{auc\_panel\_minimo}: 0,9656.
    \item \texttt{auc\_firma\_completa}: 0,9703 (AUC con los 1174 genes).
    \item \texttt{auc\_curva\_maxima}: 0,9744 (máximo de la curva, en un
    panel intermedio de 50 genes).
    \item \texttt{ihc\_recuperados}: diccionario con conteos por linaje
    (\texttt{Escamoso}: 12, \texttt{Adenocarcinoma}: 6).
    \item \texttt{ihc\_total}: diccionario con conteos totales del panel
    IHC de la OMS (\texttt{Escamoso}: 12, \texttt{Adenocarcinoma}: 8).
    \item \texttt{top10}: lista de los 10 primeros genes por $d$ media.
    \item \texttt{panel\_minimo\_genes}: lista de los 20 genes del panel
    mínimo.
  \end{itemize}

  \item[\texttt{PANEL\_MINIMO\_CURVA.csv}] Una fila por tamaño de panel
  ensayado. Columnas: \texttt{N\_Genes}, \texttt{AUC\_Media},
  \texttt{Bal\_Acc\_Media}, \texttt{AUC\_Minima}. Es la fuente de la
  figura~\ref{fig:panel_minimo}.
\end{description}

\section*{Análisis biológico complementario}
\phantomsection
\addcontentsline{toc}{section}{Análisis biológico complementario}

\begin{description}
  \item[\texttt{COMPOSICION\_VS\_BIOLOGIA.csv}] Una fila por cohorte con
  la correlación entre el score del clasificador y el contenido residual
  de pulmón sano dentro de los tumores. Columnas: \texttt{Cohorte},
  \texttt{n\_tumores}, \texttt{n\_sanos}, \texttt{Rho\_TODAS\_vs\_PulmonNormal},
  \texttt{p\_TODAS}, \texttt{Rho\_SOLO\_TUMORES\_vs\_PulmonNormal},
  \texttt{p\_SOLO\_TUMORES}, \texttt{Rho\_SOLO\_TUMORES\_vs\_Proliferacion}.
  Alimenta la figura~\ref{fig:composicion} y la sección biológica del
  capítulo~\ref{cap:resultados}.

  \item[\texttt{FALACIA\_FOLDS\_COMPARACION.csv}] Concordancia de signos
  entre folds LODO, para descartar la falacia de particiones aleatorias.
  Alimenta la figura~\ref{fig:concordancia}.

  \item[\texttt{SUBTIPO\_CASOS\_DIFICILES.csv}] Muestras cuyo score de
  probabilidad ADC/SQC queda cerca de 0,5 (candidatas a casos ambiguos).
  Base para la sección de casos difíciles del capítulo de discusión.

  \item[\texttt{SUBTIPO\_PROBS\_AMBIGUAS.csv}] Probabilidades por muestra
  y por modelo (LR-L1, RF, SVM) para las muestras seleccionadas como
  ambiguas.
\end{description}

\section*{Ficheros generados por cohorte}
\phantomsection
\addcontentsline{toc}{section}{Ficheros generados por cohorte}

Bajo cada carpeta \texttt{TFM\_GSE*/}, el pipeline persiste los
siguientes artefactos:

\begin{description}
  \item[\texttt{metadata\_procesada.csv}] Metadatos por muestra ya
  curados por el LLM. Columnas mínimas: \texttt{geo\_accession},
  \texttt{grupo\_analisis} (sano / enfermo / sin\_clasificar), y otros
  campos originales de GEO.

  \item[\texttt{expresion.csv}] Matriz de expresión gen $\times$ muestra
  ya normalizada (log$_2$ + cuantiles).

  \item[\texttt{resultados\_completos.csv}] Análisis diferencial: una
  fila por gen con \texttt{GENE\_SYMBOL}, \texttt{LogFC},
  \texttt{pvalue}, \texttt{adj\_pvalue}.

  \item[\texttt{resultados\_ml.csv}] Rendimiento local de los tres
  modelos (LR-L1, RF, SVM) sobre esa cohorte con validación cruzada
  interna. Columnas: \texttt{modelo}, \texttt{accuracy}, \texttt{auc}.

  \item[\texttt{top20\_genes\_ml\_GSE\ldots.csv}] Los 20 genes con mayor
  coeficiente absoluto en LR-L1 y su importancia en Random Forest.

  \item[\texttt{informe\_biologico.txt}] Interpretación narrativa
  automática de los DEG en función de los paneles biológicos conocidos.

  \item[\texttt{pca\_plot.png}, \texttt{volcano\_plot.png},
  \texttt{heatmap\_final.png}] Figuras exploratorias por cohorte.
\end{description}

\section*{Convenciones de tipos y unidades}
\phantomsection
\addcontentsline{toc}{section}{Convenciones de tipos y unidades}

\begin{itemize}
  \item Métricas de clasificación (AUC, balanced accuracy, sensibilidad,
  especificidad): floats en el rango $[0, 1]$, con 4 decimales de
  precisión en los CSV; se redondean a 3 decimales en la interfaz y la
  memoria.
  \item Correlaciones $\rho$: floats en el rango $[-1, 1]$; el signo es
  significativo (negativo indica relación inversa) y se preserva en
  todas las representaciones.
  \item Conteos: enteros no negativos.
  \item Símbolos génicos: nombres HUGO oficiales en mayúsculas.
  \item Cohortes: identificadores GEO Series (\texttt{GSE\ldots}).
\end{itemize}
